This section presents the proposed Causal Multiscale Wavelet Spillover Learning (CMWSL) framework. CMWSL is designed around a simple but demanding principle: any improvement over a strong volatility benchmark should be attributable to a clearly defined source of information rather than to opaque modelling complexity. For that reason, the framework is organized as a layered system. A persistence layer preserves the HAR benchmark structure. A multiscale layer extracts causal wavelet summaries from volatility-related input series. A spillover layer adds peer-index frequency-domain information. Finally, a warning layer translates the same causal feature space into high-volatility event probabilities. The full system is therefore designed for attribution as much as for prediction.

\subsection{Problem Setup}

CMWSL addresses two linked tasks: stock-index volatility forecasting as a regression problem and risk early warning as a classification problem. Let \(y_{i,t}\) denote the observed daily volatility proxy for index \(i\) on day \(t\). For a forecast horizon \(h\), the forward-looking target is defined as
\[
y_{i,t}^{(h)}=\frac{1}{h}\sum_{j=1}^{h} y_{i,t+j},
\]
so that the model predicts future average volatility rather than an in-sample quantity. This choice makes the regression task directly relevant to market-risk monitoring, where the object of interest is the volatility expected over the next few trading days rather than the volatility realized today.

Table~\ref{tab:notation} summarizes the main notation used in the paper.

\begin{table}[H]
\caption{Core notation used in the methodological formulation.}
\label{tab:notation}
\centering
\small
\setlength{\tabcolsep}{4pt}
\renewcommand{\arraystretch}{1.12}
\begin{tabularx}{\textwidth}{>{\raggedright\arraybackslash}p{0.15\textwidth}>{\raggedright\arraybackslash}p{0.18\textwidth}X}
\toprule
Symbol & Category & Meaning \\
\midrule
\(i\) & Cross-sectional index & Market index identifier (\(\mathrm{SPX}\), \(\mathrm{NDQ}\), or \(\mathrm{DJI}\)) \\
\(t\) & Time index & Forecast origin on trading day \(t\) \\
\(h\) & Forecast horizon & Prediction horizon in trading days (\(1\), \(5\), or \(10\)) \\
\(y_{i,t}^{(h)}\) & Forecast target & Future average volatility over the next \(h\) trading days \\
\(\bar y_{i,t}^{(h)}\) & Transformed target & Log-transformed target used for model fitting \\
\(x_{i,t}^{HAR}\) & Feature block & Daily, weekly, and monthly persistence features \\
\(x_{i,t}^{raw}\) & Feature block & Technical, implied-volatility, and macro-financial predictors \\
\(x_{i,t}^{wav}\) & Feature block & Causal wavelet summaries from volatility and risk series \\
\(\tilde y_{i,t}^{(h)}\) & Calibrated forecast & Post-floor volatility prediction used for evaluation and warning linkage \\
\(\tau_{i,t}^{(h)}\) & Warning threshold & Rolling in-window 90th percentile for defining high-risk states \\
\(\hat p_{i,t}^{(h)}\) & Warning score & Predicted high-risk probability from the logistic warning layer \\
\bottomrule
\end{tabularx}
\end{table}


For estimation, the target is transformed as
\[
\bar y_{i,t}^{(h)}=\log\left(1+y_{i,t}^{(h)}+\varepsilon\right),
\]
where \(\varepsilon>0\) is a small numerical constant. The back-transformed forecast is
\[
\hat y_{i,t}^{(h)}=\exp\left(\widehat{\bar y}_{i,t}^{(h)}\right)-1-\varepsilon.
\]
The transformation preserves positivity and reduces the influence of extreme observations, which is useful when tree-based models are optimized under rolling samples containing both tranquil and crisis periods.

The main empirical target is the Rogers--Satchell volatility proxy,
\[
y_{i,t}^{RS}
=
\log\!\left(\frac{H_{i,t}}{C_{i,t}}\right)\log\!\left(\frac{H_{i,t}}{O_{i,t}}\right)
+
\log\!\left(\frac{L_{i,t}}{C_{i,t}}\right)\log\!\left(\frac{L_{i,t}}{O_{i,t}}\right),
\]
where \(O_{i,t}\), \(H_{i,t}\), \(L_{i,t}\), and \(C_{i,t}\) denote open, high, low, and close prices, respectively. The robustness target is the Parkinson estimator,
\[
y_{i,t}^{P}=\frac{1}{4\log 2}\left[\log\!\left(\frac{H_{i,t}}{L_{i,t}}\right)\right]^2.
\]
Both targets exploit OHLC information more efficiently than pure close-to-close variance, while remaining compatible with a positive-target forecasting framework.

\subsection{Within-Index Causal Multiscale Representation}

The key modelling premise is that stock-index volatility is generated by heterogeneous components evolving at different temporal scales. A single lag structure is therefore unlikely to capture the whole signal. CMWSL addresses this by applying a causal non-decimated stationary wavelet transform (SWT) to selected input series \cite{percival2000wavelet}. The SWT (also known as the \`{a} trous or undecimated wavelet transform) yields translation-invariant coefficient sequences $\{d_t^{(j)}\}$ and $\{a_t^{(J)}\}$ at each scale $j$. Crucially, the causal implementation adopted here computes all coefficients using only observations $\{s_\tau:\tau\leq t\}$, achieved by applying one-sided filter banks so that no future information enters any coefficient. The non-decimated structure preserves the original sampling rate at every scale, allowing direct alignment between wavelet features and future volatility targets without interpolation or re-indexing---a necessary property for a risk model that must be operationally deployable in real time.

The implementation uses a Symlet-4 basis with \(J=3\) decomposition levels. Wavelet features are computed only after a minimum history of 128 observations, using at most the most recent \(N=256\) observations at each forecast origin. Symlet-4 is chosen because it offers near-symmetry, limited phase distortion, and a support length short enough for causal rolling application while retaining enough vanishing moments to capture smooth local structure \cite{percival2000wavelet,souropanis2023wavelet}. Three decomposition levels are used because they correspond approximately to 2-, 4-, and 8-day scales, which map naturally onto the daily, weekly, and biweekly persistence structure underlying HAR-type volatility dynamics \cite{corsi2008har,andersen2003rvol}.

For a univariate input series \(s_t\), the multi-resolution representation is
\[
s_t=a_t^{(J)}+\sum_{j=1}^{J} d_t^{(j)},
\]
where \(a_t^{(J)}\) is the smooth approximation at level \(J\) and \(d_t^{(j)}\) is the detail component at scale \(j\). CMWSL does not feed raw wavelet coefficients directly into the learner. Instead, it constructs a compact summary representation from the recent history of each coefficient sequence. For scale \(j\) and rolling window \(w\), the main summary statistics are
\[
\mu_{t,w}^{(j)}=\frac{1}{w}\sum_{\ell=0}^{w-1} d_{t-\ell}^{(j)},
\qquad
\sigma_{t,w}^{(j)}=
\sqrt{\frac{1}{w}\sum_{\ell=0}^{w-1}\left(d_{t-\ell}^{(j)}-\mu_{t,w}^{(j)}\right)^2},
\]
and
\[
E_{t,w}^{(j)}=\frac{1}{w}\sum_{\ell=0}^{w-1}\left(d_{t-\ell}^{(j)}\right)^2.
\]
These statistics summarize local level, dispersion, and energy at each scale. In practice, this is preferable to using raw coefficients because it produces a stable, low-dimensional, and interpretable feature block suitable for rolling estimation.

Let the full predictor vector for index \(i\) at time \(t\) be
\[
x_{i,t}=
\left[
x_{i,t}^{HAR},\,
x_{i,t}^{tech},\,
x_{i,t}^{macro},\,
x_{i,t}^{wav}
\right],
\]
where \(x_{i,t}^{HAR}\) contains persistence terms, \(x_{i,t}^{tech}\) contains return-transition and technical-indicator features, \(x_{i,t}^{macro}\) contains public exogenous risk variables, and \(x_{i,t}^{wav}\) contains the SWT summary block. This flat, interpretable feature vector is the core within-index feature space used by \textit{wavelet\_lightgbm}.

\subsection{Cross-Index Wavelet Spillover Features}
\label{subsec:spillover_features}

The within-index wavelet representation is informative, but it treats each index as if its volatility dynamics were self-contained. That assumption is often too strong in market-risk applications, especially during system-wide stress. CMWSL therefore extends the within-index representation with a cross-index spillover layer built from peer-market wavelet energy.

Let \(\mathcal{I}\) be the set of indices under study, \(\mathcal{I}=\{\text{DJIA},\text{NDQ},\text{SPX}\}\) in the main experiment. For each source index \(k\neq i\), let \(d_{k,t}^{(j)}\) denote the scale-\(j\) wavelet detail coefficient extracted from source \(k\)'s time series. The spillover feature passed to the model for target index \(i\) is defined as
\[
\phi_{k\to i,t,w}^{(j)}=
E_{k,t,w}^{(j)}
=
\frac{1}{w}\sum_{\ell=0}^{w-1}\left(d_{k,t-\ell}^{(j)}\right)^2,
\]
where \(j\in\{2,3\}\) and \(w\) belongs to the same set of rolling windows used for within-index wavelet summaries. The focus on levels \(j=2\) and \(j=3\) is deliberate. These correspond roughly to weekly and biweekly-to-monthly fluctuations, the scales at which slower cross-market transmission is more plausible than at the highest-frequency noise level \cite{son2023gnnspillover,souropanis2023wavelet}.

The spillover block is constructed from three source series for each peer index: the realized volatility proxy \(y_{k,t}^{RS}\), the absolute close-to-close return \(|r_{k,t}^{cc}|\), and the index-specific implied-volatility proxy. After strict causal alignment, meaning that only information observed by time \(t\) is used to forecast \(y_{i,t}^{(h)}\), the empirical implementation appends 126 distinct spillover columns to the panel. An active-column filter is applied in each rolling training window, retaining only columns with sufficient non-missing coverage. This avoids self-reference artifacts and prevents the effective predictor dimension from being inflated by structurally empty variables.

The extended predictor vector for the spillover model is therefore
\[
x_{i,t}^{spill}=
\left[
x_{i,t}^{HAR},\,
x_{i,t}^{tech},\,
x_{i,t}^{macro},\,
x_{i,t}^{wav},\,
\bm{\phi}_{k\to i,t}
\right],
\]
where \(\bm{\phi}_{k\to i,t}\) collects all active peer-index spillover features. The critical identification feature of the empirical design is that \textit{spillover\_lightgbm} is otherwise identical to \textit{wavelet\_lightgbm}; any difference in out-of-sample performance is therefore attributable to the added spillover block rather than to a different learner or different hyperparameters.

\subsection{Forecasting Layer}

The forecasting layer combines a strong persistence benchmark with nonlinear interaction learning. The HAR block is defined as
\[
x_{i,t}^{HAR}=
\left(
y_{i,t},
\frac{1}{5}\sum_{\ell=0}^{4}y_{i,t-\ell},
\frac{1}{22}\sum_{\ell=0}^{21}y_{i,t-\ell}
\right),
\]
which preserves the benchmark view that volatility persistence is heterogeneous across daily, weekly, and monthly horizons \cite{corsi2008har}. Short-run transition features are defined from OHLC prices as
\[
r_{i,t}^{cc}=\log C_{i,t}-\log C_{i,t-1},\qquad
r_{i,t}^{oc}=\log C_{i,t}-\log O_{i,t},
\]
\[
g_{i,t}^{co}=\log O_{i,t}-\log C_{i,t-1}.
\]
These terms capture close-to-close movement, intraday dynamics, and overnight gaps, all of which are relevant for daily index-volatility formation.

At the forecasting stage, a horizon-specific LightGBM model is fitted on the log-transformed training target over the rolling window and used to generate a one-step-ahead out-of-sample forecast. The spillover variant uses the same procedure with the spillover-augmented feature vector.

LightGBM \cite{ke2017lightgbm} is used as the main nonlinear learner for three reasons. First, the rolling protocol requires hundreds of repeated refits, so computationally light models are preferred to heavyweight sequence models. Second, gradient-boosted trees handle mixed feature blocks and nonlinear interactions without requiring feature rescaling assumptions. Third, split-based feature attribution and SHAP analysis remain straightforward, which is important because the contribution of the paper depends on identifying what kind of information drives forecast gains rather than only on reporting a rank order of models.

Because volatility targets are strictly nonnegative and QLIKE is sensitive to unrealistically small predictions, the framework applies a horizon-specific floor calibration:
\[
\tilde y_{i,t}^{(h)}=
\max\left(\hat y_{i,t}^{(h)},q_{\alpha,h}^{\text{train}}\right),
\]
where \(q_{\alpha,h}^{\text{train}}\) is a low empirical quantile estimated from the current training window. In the implementation, the floor quantile is \(0.05\) for \(h=1\) and \(0.01\) for \(h=5\) and \(h=10\). This calibration improves numerical stability without using any future information.

\begin{figure}[!t]
\centering
\includegraphics[width=0.96\textwidth]{figures/Method.png}
\caption{Overall architecture of the proposed CMWSL stock-index volatility forecasting and risk-warning framework.}
\label{fig:model_architecture}
\end{figure}

Figure~\ref{fig:model_architecture} summarizes the complete data-to-feature-to-prediction pipeline. An expanded module-level diagram is provided in the appendix for readers who want a more granular view of the persistence, raw-risk, and multiscale branches.

\subsection{Warning Layer}

The risk-warning task converts future volatility realizations into an operational high-risk label. For each rolling training window, the warning threshold \(\tau_{i,t}^{(h)}\) is defined as the 90th percentile of in-window future volatility targets. The binary event label is
\[
z_{i,t}^{(h)}=\mathbb{I}\!\left(y_{i,t}^{(h)}>\tau_{i,t}^{(h)}\right).
\]
This design makes the event definition adaptive to level differences across indices and to changing market regimes.

The main warning specification uses a logistic model estimated on current market-risk features:
\[
x_{i,t}^{warn}=\left[x_{i,t}^{raw},\,s_{i,t}^{stress}\right],
\]
where \(s_{i,t}^{stress}\) collects short-window stress summaries derived from the public predictor set. The probability of a future high-volatility event is modeled as
\[
\Pr\!\left(z_{i,t}^{(h)}=1\mid x_{i,t}^{warn}\right)
=
\frac{1}{1+\exp\left[-\left(\beta_0+\beta^\top x_{i,t}^{warn}\right)\right]}.
\]
The final warning decision is then
\[
\hat z_{i,t}^{(h)}=
\mathbb{I}\!\left(\hat p_{i,t}^{(h)}\ge c_{i,h}^{\star}\right),
\]
with the threshold chosen by
\[
c_{i,h}^{\star}
=
\arg\max_{c\in[0,1]}
F_{\beta}\!\left(y^{train},\mathbb{I}(\hat p^{train}\ge c)\right),
\qquad \beta=2.
\]
Using \(\beta=2\) reflects the asymmetry of practical warning problems, where missed high-volatility events are usually more costly than additional false alarms.

Finally, the main regression loss used throughout the paper is QLIKE,
\[
L_{QLIKE}(y,\hat y)=\log(\hat y)+\frac{y}{\hat y},
\]
which is standard for volatility forecasting because it rewards both accuracy and calibration on the positive target scale \cite{patton2011proxy}. This is particularly appropriate here because the empirical objective is not merely to reduce squared error, but to produce stable positive volatility forecasts that remain usable inside a market-risk monitoring system.

\begin{table}[!t]
\caption{Leakage-free pseudo-code of the rolling multiscale forecasting and warning procedure.}
\label{tab:pseudocode}
\centering
\small
\setlength{\tabcolsep}{5pt}
\renewcommand{\arraystretch}{1.12}
\begin{tabularx}{\textwidth}{>{\raggedright\arraybackslash}p{0.08\textwidth}X}
\toprule
Step & Operation \\
\midrule
1 & For each index \(i\) and horizon \(h\), define a rolling training window of 1260 trading days and a one-step-ahead out-of-sample forecast origin. \\
2 & Construct the OHLC-based volatility target \(y_{i,t}^{(h)}\), persistence terms, technical indicators, implied-volatility proxies, and macro-financial predictors using only data available up to time \(t\). \\
3 & Apply the causal non-decimated wavelet transform to selected high-value series and compute multiscale summaries such as recent coefficients, short-window means, standard deviations, and local energy terms. \\
4 & Form the feature vector \(x_{i,t}=[x_{i,t}^{HAR},x_{i,t}^{raw},x_{i,t}^{wav}]\) and estimate the horizon-specific boosting model \(f_h(\cdot)\) on the transformed training target. \\
5 & Generate the raw forecast \(\hat y_{i,t}^{(h)}\) and apply the horizon-specific lower-tail calibration safeguard to obtain the final regression forecast \(\tilde y_{i,t}^{(h)}\). \\
6 & Compute the rolling high-risk threshold \(\tau_{i,t}^{(h)}\), estimate the logistic warning model on the training window, and choose the decision threshold that maximizes the training-window \(F_{\beta}\) score with \(\beta=2\). \\
7 & Output the calibrated volatility forecast, the warning probability, and the warning label for the current out-of-sample date, then roll the window forward and repeat. \\
\bottomrule
\end{tabularx}
\end{table}


Table~\ref{tab:pseudocode} summarizes the leakage-free rolling implementation of CMWSL.
